\clearpage
\thispagestyle{empty}\SourceLeaf{249}
\pdfbookmark[0]{Dioptrical Investigations}{dioptrische-title}
\begin{center}
\vspace*{0.6in}
{\Large DIOPTRICAL INVESTIGATIONS\par}
\vspace{0.6in}
BY\par
\vspace{0.5in}
{\large CARL FRIEDRICH GAUSS\par}
\vspace{0.4in}
PRESENTED TO THE ROYAL SOCIETY MDCCCXL DECEMBER X.\par
\vfill
\rule{0.85\textwidth}{0.3pt}\par
Transactions of the Royal Society of Sciences in Göttingen, Volume I.\par
Göttingen \textsc{mdcccxliii}.\par
\rule{0.85\textwidth}{0.3pt}
\end{center}
\clearpage
\SourceLeaf{251}
\fancyhead[L]{\small Dioptrical Investigations}
\renewcommand{\Article}[1]{\par\medskip\Needspace{4\baselineskip}\pdfbookmark[1]{Art. #1}{dioptrische-art-#1}\centerline{\textbf{#1.}}\nobreak\smallskip}
\pdfbookmark[0]{Dioptrical Investigations}{dioptrische-body}
\begin{center}
DIOPTRICAL INVESTIGATIONS.
\end{center}
\medskip
The study of the path taken through glass lenses by light rays inclined very little to their common axis, and of the phenomena depending on it, yields very elegant results. The work of \textsc{Cotes}, \textsc{Euler}, \textsc{Lagrange} and \textsc{Möbius} might seem to have exhausted these results, yet much remains to be desired. An essential shortcoming of the propositions established by those mathematicians is that they neglect the thickness of the lenses. This gives them an inaccurate and unnatural character that greatly diminishes their value. The initial neglect of lens thickness is very useful, indeed necessary, for some other dioptrical investigations, particularly those taking account of the so-called aberration due to the spherical shape of the lens surfaces, since it gives simpler and more flexible rules for estimates and first approximations. Without denying this, one will nevertheless gladly be spared such a sacrifice of all precision wherever this can be done with no loss, or no substantial loss, in the simplicity of the results. Even in the initial definitions of dioptrics we already encounter, to some extent, a lack of precision that offends the mathematical sense. The concepts of the axis and focal point of a lens are precisely established; the same is not true of
\SourceLeaf{252}focal length, which most writers define as the distance of the focal point of the lens from its centre. From the outset they either tacitly assume, or expressly state, that the thickness of the lens is to be treated as infinitely small. For actual lenses, therefore, the focal length retains an indeterminacy of the order of the lens thickness. Where greater accuracy is attempted, this distance is measured sometimes from the lens surface nearest the focal point, sometimes from its so-called optical centre, and sometimes from the point midway between its front and rear surfaces. Different again from all these determinations is the value that must underlie the comparison of the size of the image of an infinitely distant object with the apparent size of that object; this last determination is in fact the only appropriate one.

I have therefore not thought it superfluous to devote a few pages to these investigations, elementary as they are in themselves, chiefly to show that the thickness of the lenses can be taken into account in the elegant propositions mentioned above without any loss of simplicity. Only the restriction to rays inclined infinitely little to the axis is to be retained here, or the aberration due to the spherical shape is to be set aside.

\Article{1}
The position of every point occurring in this investigation is determined by rectangular coordinates $x,y,z$. It is assumed that the centres of the various refracting surfaces lie on the $x$-axis, and only light rays making a very small angle with this axis are considered: the coordinates $x$, with an entirely arbitrary origin, are taken to increase in the direction of the light rays.

We first consider the effect of a single refraction on the path of a light ray. Let the ratio of refraction on passing from the first medium into the second be as $\frac1n$ to $\frac1{n'}$, or as $n'$ to $n$. We denote by $M$ the centre of the spherical surface separating the two media, and by $N$ the intersection of this surface with the first coordinate axis; the same letters are also to denote the corresponding
\SourceLeaf{253}values of $x$ for these points, a convention that will likewise be followed below for other points on the first coordinate axis. Further, let $r=M-N$, so that $r$ is the radius of the separating surface, positive or negative according as the first medium lies on the convex or the concave side; let $P$ be the point where the light ray meets the separating surface, and $\theta$ the (acute) angle between $MP$ and the $x$-axis.

The straight line described by a light ray before refraction is determined by two equations, which we put in the following form:
\[
\begin{aligned}
y&=\frac\beta n(x-N)+b\\
z&=\frac\gamma n(x-N)+c
\end{aligned}
\]
and similarly let the equations of the straight line described by the same light ray after refraction be
\[
\begin{aligned}
y&=\frac{\beta'}{n'}(x-N)+b'\\
z&=\frac{\gamma'}{n'}(x-N)+c'
\end{aligned}
\]
We must therefore determine how the four quantities $\beta',\gamma',b',c'$ depend on $\beta,\gamma,b,c$. For the point $P$ we have
\[
x=N+r(1-\cos\theta)
\]
and since both the first and the second equations hold for this point,
\[
\frac\beta n.r(1-\cos\theta)+b=\frac{\beta'}{n'}.r(1-\cos\theta)+b'
\]
Consequently, since $\beta,\beta',\theta$ count as infinitely small quantities of first order, we have, accurate apart from quantities of third order,
\[
b'=b\quad .\quad .\quad .\quad .\quad .\quad .\quad .\quad .\quad .\quad .\quad .\qquad(1)
\]
and likewise
\[
c'=c\quad .\quad .\quad .\quad .\quad .\quad .\quad .\quad .\quad .\quad .\quad .\qquad(1)
\]

Let a plane through $M$ perpendicular to the $x$-axis meet the first path of the light ray (extended if necessary) at $Q$ and the second at $Q'$. Since $PQ'$ lies in one plane with $PQ$ and $PM$, the points $M,Q,Q'$ lie on one straight line. If $\lambda,\lambda'$ denote the angles this straight line makes with $PQ,PQ'$, then evidently
\SourceLeaf{254}$MQ.\sin\lambda,MQ'.\sin\lambda'$ equal the products of the radius of the spherical surface, taken positively, with the sines of the angle of incidence and the angle of refraction, respectively, and are therefore proportional to the numbers $n',n$. Hence
\[
MQ'=\frac{n.MQ.\sin\lambda}{n'\sin\lambda'}
\]
Now for the point $Q$ we have
\[
\begin{aligned}
y&=b+\frac{\beta r}n\\
z&=c+\frac{\gamma r}n
\end{aligned}
\]
whereas for the point $Q'$,
\[
\begin{aligned}
y&=b'+\frac{\beta'r}{n'}\\
z&=c'+\frac{\gamma'r}{n'}
\end{aligned}
\]
and since the latter two coordinates are to the former two as $MQ'$ to $MQ$, we obtain
\[
\begin{aligned}
b'+\frac{\beta'r}{n'}&=\frac{n\sin\lambda}{n'\sin\lambda'}.\left(b+\frac{\beta r}n\right)\\
c'+\frac{\gamma'r}{n'}&=\frac{n\sin\lambda}{n'\sin\lambda'}.\left(c+\frac{\gamma r}n\right)
\end{aligned}
\]
or
\[
\begin{aligned}
\beta'&=\frac{nb+\beta r}{r}.\frac{\sin\lambda}{\sin\lambda'}-\frac{n'b'}r\\
\gamma'&=\frac{nc+\gamma r}{r}\raisebox{0.55em}{.}\frac{\sin\lambda}{\sin\lambda'}-\frac{n'c'}r
\end{aligned}
\]
These expressions hold exactly; but since $\lambda,\lambda'$ differ from a right angle by quantities of first order, and hence their sines differ from unity by quantities of second order, we have, accurate apart from quantities of third order,
\[
\left.\begin{aligned}
\beta'&=\beta-\frac{n'-n}r.b=\beta+\frac{n'-n}{N-M}.b\\
\gamma'&=\gamma-\frac{n'-n}r.c=\gamma+\frac{n'-n}{N-M}.c
\end{aligned}\right\}\quad .\quad .\quad .\quad .\quad .\quad .\qquad(2)
\]
These equations (1), (2) give the solution of our problem.

It deserves mention that the same formulas can also be applied directly to a reflected ray if one merely substitutes $-n$ for $n'$, and that this procedure also makes it very easy to extend all the following investigations to the case where one or more reflections take the place of the refractions.

\SourceLeaf{255}
\Article{2}
To solve the general problem of determining the path of a light ray after $(\mu+1)$ successive refractions, we shall use the following notation.

$N^0,N',N''\, .\, .\, .\, .\, .\, N^{(\mu)}$ are the points where the refracting surfaces meet the $x$-axis.

$M^0,M',M''\, .\, .\, .\, .\, M^{(\mu)}$ are the centres, on this axis, of the refracting surfaces

$n':n^0,\ n'':n',\ n''':n''\, .\, .\, .\, .\, .\, n^{(\mu+1)}:n^{(\mu)}$ are the ratios of refraction on passing from the first medium (before $N^0$) into the second (between $N^0$ and $N'$), from the second into the third, and so on. In the emission theory, therefore, the numbers $n^0,n',n''$ and so on are directly proportional to the speeds of propagation of light in the individual media, and in the wave theory inversely proportional; if the last medium is the same as the first, then $n^{(\mu+1)}=n^0$.

Let the equations for the path of the light ray before the first refraction be
\[
\begin{aligned}
y&=\frac{\beta^0}{n^0}(x-N^0)+b^0\\
z&=\frac{\gamma^0}{n^0}(x-N^0)+c^0
\end{aligned}
\]
and those for the path after the first refraction
\[
\begin{aligned}
y&=\frac{\beta'}{n'}(x-N^0)+b^0\\
z&=\frac{\gamma'}{n'}(x-N^0)+c^0
\end{aligned}
\]
or, with reference to $N^0$ replaced by reference to $N'$,
\[
\begin{aligned}
y&=\frac{\beta'}{n'}(x-N')+b'\\
z&=\frac{\gamma'}{n'}(x-N')+c'
\end{aligned}
\]
and likewise the equations for the path after the second refraction
\[
\begin{aligned}
y&=\frac{\beta''}{n''}(x-N')+b'\\
z&=\frac{\gamma''}{n''}(x-N')+c'
\end{aligned}
\]
or
\[
\begin{aligned}
y&=\frac{\beta''}{n''}(x-N'')+b''\\
z&=\frac{\gamma''}{n''}(x-N'')+c''
\end{aligned}
\]
\SourceLeaf{256}and so on. Thus, if we denote the last terms in the sequences of $\beta,\gamma,n,N,b,c$, namely $\beta^{(\mu+1)},\gamma^{(\mu+1)},n^{(\mu+1)},N^{(\mu)},b^{(\mu)},c^{(\mu)}$, by $\beta^*,\gamma^*,n^*,N^*,b^*,c^*$ to distinguish them as the last terms, the equations for the final path of the light ray are
\[
\begin{aligned}
y&=\frac{\beta^*}{n^*}(x-N^*)+b^*\\
z&=\frac{\gamma^*}{n^*}(x-N^*)+c^*
\end{aligned}
\]
Finally, for brevity we set
\[
\left.\begin{aligned}
\frac{N'-N^0}{n'}&=t',&\frac{N''-N'}{n''}&=t'',&\frac{N'''-N''}{n'''}&=t'''\ \text{etc.}\\
\frac{n'-n^0}{N^0-M^0}&=u^0,&\frac{n''-n'}{N'-M'}&=u',&\frac{n'''-n''}{N''-M''}&=u''\ \text{etc.}
\end{aligned}\right\}\quad .\quad .\qquad (3)
\]
and, by analogy, for the last terms in these sequences,
\[
t^{(\mu)}=t^*,\qquad u^{(\mu)}=u^*
\]
It follows from the preceding article that
\[
\begin{aligned}
\beta'&=\beta^0+u^0b^0\\
b'&=b^0+t'\beta'\\
\beta''&=\beta'+u'b'\\
b''&=b'+t''\beta''\\
\beta'''&=\beta''+u''b''\\
b'''&=b''+t'''\beta'''
\end{aligned}
\]
and so on. We see from this that $b^*,\beta^*$ are determined linearly by $b^0$ and $\beta^0$, and that, if we set
\[
\left.\begin{aligned}
b^*&=gb^0+h\beta^0\\
\beta^*&=kb^0+l\beta^0
\end{aligned}\right\}\quad .\quad .\quad .\quad .\quad .\quad .\quad .\quad .\quad .\qquad (4)
\]
then in the notation introduced by \textsc{Euler} (\textit{Comment. Nov. Acad. Petropol.} T. IX) we have
\[
\left.\begin{aligned}
g&=(u^0,t',u',t'',u''\ .\ .\ .\ .\ .\ t^*)\\
h&=(t',u',t'',u''\ .\ .\ .\ .\ .\ t^*)\\
k&=(u^0,t',u',t'',u''\ .\ .\ .\ .\ .\ .\ u^*)\\
l&=(t',u',t'',u''\ .\ .\ .\ .\ .\ .\ u^*)
\end{aligned}\right\}\quad .\quad .\quad .\quad .\quad .\quad .\quad .\qquad (5)
\]
\SourceLeaf{257}As is well known, this notation means that if, from a given sequence of quantities $a,a',a'',a'''$ and so on, another sequence, $A,A',A'',A'''$ and so on, is formed by the following algorithm,
\[
A=a,\quad A'=a'A+1,\quad A''=a''A'+A,\quad A'''=a'''A''+A'\quad\text{etc.}
\]
we write
\[
A=(a),\quad A'=(a,a'),\quad A''=(a,a',a''),\quad A'''=(a,a',a'',a''')\quad\text{etc.}
\]
It is also evident that in the equations for the third coordinate $z$ the constants for the final path are derived from those for the first in exactly the same way as in the equations for $y$, or that we have
\[
\left.\begin{aligned}
c^*&=gc^0+h\gamma^0\\
\gamma^*&=kc^0+l\gamma^0
\end{aligned}\right\}\quad .\quad .\quad .\quad .\quad .\quad .\quad .\quad .\quad .\qquad (4)
\]
Equations (3), (5), (4) contain the complete solution of our problem.

\Article{3}
In the work cited, \textsc{Euler} developed the principal relations concerning the algorithm just mentioned; only two will be recalled here.

First, we always have
\[
(a,a',a''\ldots a^{(\lambda)})(a',a''\ldots a^{(\lambda+1)})-(a,a',a''\ldots a^{(\lambda+1)})(a',a''\ldots a^{(\lambda)})=\pm1
\]
where the upper or lower sign applies according as the total number of elements $a,a',a''\ldots a^{(\lambda+1)}$, that is, the number $\lambda+2$, is odd or even.

Second, the order of the elements may be reversed; namely,
\[
(a,a',a''\ldots a^{(\lambda)})=(a^{(\lambda)},\ldots a'',a'\ a)
\]
Applying the first of these propositions to the quantities $g,h,k,l$ gives
\[
gl-hk=1
\]
Equations (4) can therefore also be written as follows:
\SourceLeaf{258}
\[
\begin{aligned}
b^0&=lb^*-h\beta^*\\
\beta^0&=-kb^*+g\beta^*\\
c^0&=lc^*-h\gamma^*\\
\gamma^0&=-kc^*+g\gamma^*
\end{aligned}
\]

\Article{4}
Let $P$ be a given point on the straight line representing the first path of the light ray (extended if necessary), and let $\xi,\eta,\zeta$ be its coordinates. We therefore have
\[
n^0\eta=\beta^0(\xi-N^0)+n^0b^0
\]
or, substituting for $\beta^0,b^0$ the expressions given at the end of the preceding article,
\[
n^0\eta=(g\beta^*-kb^*)(\xi-N^0)-n^0(h\beta^*-lb^*)
\]
and hence
\[
b^*=\frac{n^0\eta+(n^0h-g(\xi-N^0))\beta^*}{n^0l-k(\xi-N^0)}
\]
If this value is substituted into the first equation for the path of the light ray after the last refraction, namely,
\[
y=\frac{\beta^*}{n^*}(x-N^*)+b^*
\]
and for brevity we write
\[
\begin{aligned}
N^*-\frac{n^0h-g(\xi-N^0)}{n^0l-k(\xi-N^0)}.n^*&=\xi^*\\
\frac{n^0\eta}{n^0l-k(\xi-N^0)}&=\eta^*
\end{aligned}
\]
then this equation becomes
\[
y=\eta^*+\frac{\beta^*}{n^*}(x-\xi^*)
\]
In exactly the same way, if we also write
\[
\frac{n^0\zeta}{n^0l-k(\xi-N^0)}=\zeta^*
\]
the second equation for the path of the light ray after the last refraction becomes
\SourceLeaf{259}
\[
z=\zeta^*+\frac{\gamma^*}{n^*}(x-\xi^*)
\]
The point $P^*$, whose coordinates are $\xi^*,\eta^*,\zeta^*$, therefore lies on the straight line representing this final path (extended backwards if necessary). Since its coordinates are independent of $\beta^0,b^0,\gamma^0,c^0$, it is also clear that this point is the same for \textit{all} incident rays passing through $P$. The point $P$ may be regarded as an object and $P^*$ as its image; the object can be real only if $P$ lies in the first medium, or $\xi-N^0$ is negative, and likewise the image is real only if $P^*$ lies in the last medium, or $\xi^*-N^*$ is positive. In the opposite cases the object or image is merely virtual.

The points $P,P^*$ lie in one plane with the $x$-axis, at distances from it in the ratio of unity to the number $\frac{n^0}{n^0l-k(\xi-N^0)}$. The positive or negative sign of this number indicates that the points lie on the same side or on opposite sides of the axis. A system of points in one plane perpendicular to the $x$-axis may be regarded as a composite object whose composite image likewise lies in a plane perpendicular to the $x$-axis and is similar to the object, so that the linear ratio of corresponding parts is expressed by the number
\[
\frac{n^0}{n^0l-k(\xi-N^0)}=g+\frac{k}{n^*}(\xi^*-N^*)
\]
whose sign distinguishes the upright from the inverted orientation.

\Article{5}
What has been developed so far contains the entire theory of the changes that the paths of light rays undergo through refraction, and can easily be extended to the case where refractions are combined with one or more reflections; that case, however, is not to be worked out here in detail. It is nevertheless not superfluous to put the results into another form, referring them, instead of to the first and last surfaces or to the points $N^0,N^*$, to two other points $Q,Q^*$. Let
\[
\begin{aligned}
y&=\frac{\beta^0}{n^0}(x-Q)+B\\
z&=\frac{\gamma^0}{n^0}(x-Q)+C
\end{aligned}
\]
\SourceLeaf{260}be the equations for the first path, and
\[
\begin{aligned}
y&=\frac{\beta^*}{n^*}(x-Q^*)+B^*\\
z&=\frac{\gamma^*}{n^*}(x-Q^*)+C^*
\end{aligned}
\]
the equations for the final path of the light ray, and set
\[
\frac{N^0-Q}{n^0}=\theta,\qquad \frac{Q^*-N^*}{n^*}=\theta^*
\]
We therefore have
\[
\begin{aligned}
b^0&=B+\theta\beta^0,&c^0&=C+\theta\gamma^0\\
B^*&=b^*+\theta^*\beta^*,&C^*&=c^*+\theta^*\gamma^*
\end{aligned}
\]
Combining these with equations (4), it follows readily that, if we set
\[
\begin{aligned}
G&=g+\theta^*k\\
H&=h+\theta g+\theta\theta^*k+\theta^*l\\
K&=k\\
L&=l+\theta k
\end{aligned}
\]
then
\[
\begin{aligned}
B^*&=GB+H\beta^0,&C^*&=GC+H\gamma^0\\
\beta^*&=KB+L\beta^0,&\gamma^*&=KC+L\gamma^0
\end{aligned}
\]
The coefficients $G,H,K,L$, which thus take the place of $g,h,k,l$, also give the equation
\[
GL-HK=1
\]

\Article{6}
The purpose of introducing other points to which the positions of the incident and outgoing rays are referred is to express more simply how the latter position depends on the former. Two pairs of points are particularly suitable for this; they will be denoted by $E,E^*$ and $F,F^*$. The values of the quantities involved can conveniently be surveyed in tabular form.
\SourceLeaf{261}
\[
\renewcommand{\arraystretch}{1.8}
\begin{array}{l|l|l}
&\multicolumn{1}{c|}{\mathrm{I}}&\multicolumn{1}{c}{\mathrm{II}}\\\cline{2-3}
\theta&\dfrac{1-l}{k}&-\dfrac{l}{k}\\
\theta^*&\dfrac{1-g}{k}&-\dfrac{g}{k}\\
Q&E=N^0-\dfrac{n^0(1-l)}{k}&F=N^0+\dfrac{n^0l}{k}=E+\dfrac{n^0}{k}\\
Q^*&E^*=N^*+\dfrac{n^*(1-g)}{k}&F^*=N^*-\dfrac{n^*g}{k}=E^*-\dfrac{n^*}{k}\\
G&1&0\\
H&0&-\dfrac1k\\
K&k&k\\
L&1&0
\end{array}
\]
The result is therefore that, if the equations for the incident ray are put in the form
\[
\begin{aligned}
y&=\frac{\beta^0}{n^0}(x-E)+B\\
z&=\frac{\gamma^0}{n^0}(x-E)+C
\end{aligned}
\]
or in the following form (where we distinguish the constant parts from those in the first form by primes),
\[
\begin{aligned}
y&=\frac{\beta^0}{n^0}(x-F)+B'\\
z&=\frac{\gamma^0}{n^0}(x-F)+C'
\end{aligned}
\]
the equations for the outgoing ray will be
\[
\begin{aligned}
y&=\frac{\beta^0+kB}{n^*}.(x-E^*)+B\\
z&=\frac{\gamma^0+kC}{n^*}.(x-E^*)+C
\end{aligned}
\]
or
\[
\begin{aligned}
y&=\frac{kB'}{n^*}.(x-F^*)-\frac{\beta^0}k\\
z&=\frac{kC'}{n^*}.(x-F^*)-\frac{\gamma^0}k
\end{aligned}
\]

\SourceLeaf{262}\Article{7}
Using the points $E,E^*$, the dependence of the final path of the light ray on its first path can be expressed simply as follows: relative to the point $E^*$, the final path has the same position as a ray refracted only once would have relative to $E$ if there were at $E$ a refracting surface of radius $\frac{n^0-n^*}{k}$, through which the light ray would pass directly from the first medium into the last. This applies when the first and last media are different. If, however, they are the same, or $n^*=n^0$, as in refraction through one or more glass lenses, the final path has the same position relative to $E^*$ as it would have relative to $E$ after refraction through an infinitely thin lens at $E$ with focal length $-\frac{n^0}k$. In other words, the transition from the first medium into the last through several refractions may be replaced either by a single refraction or by a single lens of infinitely small thickness, according as the first and last media are different or the same. In the first case we give the refracting surface the radius $\frac{n^0-n^*}{k}$, and in the second we give the lens the focal length $-\frac{n^0}k$; we place the refracting surface or the lens at $E$, and in both cases shift the position of the outgoing ray by the distance of the point $E^*$ from $E$. The sign of the radius of the refracting surface is to be understood as above in Article 1, and the sign of the focal length as will be stated below in Article 9.

These points $E,E^*$ are sufficiently significant that they seem to deserve a name of their own: I shall call them the \textit{principal points} of the system of media, the lens, or the system of lenses to which they refer; $E$ the first, and $E^*$ the second principal point. The planes of the principal points will be understood to be the planes drawn through them perpendicular to the $x$-axis,

\Article{8}
As for the points $F,F^*$, the formulas of Article 6 show that all incident light rays passing through the point $F$ correspond to outgoing rays parallel to the axis; incident rays parallel to the axis, on the other hand, correspond to outgoing rays intersecting at the point $F^*$; for
\SourceLeaf{263}rays coming from the opposite side, these points exchange their functions. If we therefore extend the terminology used for individual lenses, $F,F^*$ may be called the focal points of the system of media or lenses to which they refer, $F$ the first and $F^*$ the second; the planes drawn through these points perpendicular to the $x$-axis may be called the focal planes. Those formulas of Article 6 also show that all rays intersecting at any other point of the first focal plane correspond to outgoing rays inclined to the axis but parallel to one another. Conversely, all incident rays parallel to one another but not to the axis correspond to outgoing rays intersecting at a point of the second focal plane different from $F^*$.

\Article{9}
With these four planes we obtain a very simple construction for the position of the outgoing ray.

Let the incident ray intersect the first focal plane at the point $(1)$ and the first principal plane at the point $(2)$; let a line parallel to $(1)(2)$ drawn through $F$ meet the first principal plane at $(3)$; let a line through $(2)$ parallel to the axis meet the second principal plane at $(4)$; finally, let a line through $(3)$ parallel to the axis meet the second focal plane at $(5)$. Then $(4)(5)$ or $(5)(4)$ gives the position of the outgoing ray. The values of the coordinates are
\[
\renewcommand{\arraystretch}{1.25}
\begin{array}{c||c|c|c}
\text{for}&x&y&z\\\hline
(1)&F&B'&C'\\
(2)&E&B&C\\
F&F&0&0\\
(3)&E&B-B'&C-C'\\
(4)&E^*&B&C\\
(5)&F^*&B-B'&C-C'
\end{array}
\]
It therefore follows from the formulas of Article 6 that the outgoing ray passes through $(4)$ and $(5)$; the former follows immediately, and the latter because
\SourceLeaf{264}
\[
\begin{aligned}
B-B'&=\frac{\beta^0}{n^0}(E-F)=-\frac{\beta^0}k\\
C-C'&=\frac{\gamma^0}{n^0}(E-F)=-\frac{\gamma^0}k
\end{aligned}
\]

In the most usual case, where $n^*=n^0$ and hence $F^*-E^*=E-F$, the construction becomes even simpler because the point $(3)$ is unnecessary; we need only determine $(1),(2),(4)$ as before and draw $(4)(5)$ parallel to $(1)E$.

If the direction of the incident ray passes through $E$, the direction of the outgoing ray always passes through $E^*$ and, in the case where $n^*=n^0$, is also parallel to the former. For simple lenses it is customary to call such a ray a principal ray.

The distances of the second focal plane from the second principal plane, and of the first principal plane from the first focal plane, or the quantities $-\frac{n^*}k,-\frac{n^0}k$, could be called the \textit{focal lengths} of the system of media, if it did not seem more appropriate to restrict this name to the case where the last medium is the same as the first and those distances are therefore equal to one another. To conform to the usual terminology, we regard the focal length as positive when the first principal point has a larger coordinate than the first focal point, so that the focal length is always expressed by $-\frac{n^0}k=-\frac{n^*}k$.

\Article{10}
As is readily seen, in the formulas given above in Article 4 for the position of the image, other points may be used in place of $N^0,N^*$, provided that $g,h,k,l$ are also replaced by the corresponding $G,H,K,L$. Choosing the principal points for this purpose, we obtain the following expressions:
\[
\begin{aligned}
\xi^*&=E^*-\frac{n^*(E-\xi)}{n^0+k(E-\xi)}\\
\eta^*&=\frac{n^0\eta}{n^0+k(E-\xi)}\\
\zeta^*&=\frac{n^0\zeta}{n^0+k(E-\xi)}
\end{aligned}
\]
The first formula can also be written in the following form
\[
\frac{n^*}{\xi^*-E^*}+\frac{n^0}{E-\xi}=-k
\]
\SourceLeaf{265}If we choose the focal points, we obtain
\[
\begin{aligned}
\xi^*&=F^*+\frac{n^0n^*}{kk(F-\xi)}\\
\eta^*&=\frac{n^0\eta}{k(F-\xi)}\\
\zeta^*&=\frac{n^0\zeta}{k(F-\xi)}
\end{aligned}
\]

Since the formulas are frequently used, they may also be given here in the form they take when the first and last media are the same and the focal length is denoted by $\varphi$.
\[
\begin{gathered}
\frac1{\xi^*-E^*}+\frac1{E-\xi}=\frac1\varphi\\
(\xi^*-F^*)(F-\xi)=\varphi\varphi\\
\eta^*=-\frac{\varphi\eta}{F-\xi}=-\frac{\eta(\xi^*-F^*)}{\varphi}\\
\zeta^*=-\frac{\varphi\zeta}{F-\xi}=-\frac{\zeta(\xi^*-F^*)}{\varphi}
\end{gathered}
\]

\Article{11}
The four auxiliary points $E,E^*,F,F^*$ cease to be applicable in the special case where $k=0$. Those points would then have to be regarded as infinitely far from the refracting surfaces. In this case one can go straight to the general formulas given above for solving the main problems. They take the following form here.

If the equations for the incident ray are written as
\[
\begin{aligned}
y&=\frac{\beta^0}{n^0}(x-N^0)+b^0\\
z&=\frac{\gamma^0}{n^0}(x-N^0)+c^0
\end{aligned}
\]
then those for the outgoing ray are
\[
\begin{aligned}
y&=\frac{l\beta^0}{n^*}.(x-N^*)+gb^0+h\beta^0\\
z&=\frac{l\gamma^0}{n^*}.(x-N^*)+gc^0+h\gamma^0
\end{aligned}
\]
If, for brevity, one sets
\[
N^*-\frac{hn^*}{l}=N^{**}
\]
or, equivalently, since $gl=1$,
\SourceLeaf{266}
\[
N^*-ghn^*=N^{**}
\]
the formulas take the still simpler form
\[
\begin{aligned}
y&=\frac{l\beta^0}{n^*}.(x-N^{**})+gb^0\\
z&=\frac{l\gamma^0}{n^*}.(x-N^{**})+gc^0
\end{aligned}
\]
For the position of the image of the point whose coordinates are $\xi,\eta,\zeta$, we obtain the coordinates
\[
\begin{aligned}
\xi^*&=N^*-ghn^*-\frac{n^*}{n^0}gg(N^0-\xi)\\
&=N^{**}-\frac{n^*}{n^0}gg(N^0-\xi)\\
\eta^*&=g\eta\\
\zeta^*&=g\zeta
\end{aligned}
\]
It follows that the point on the $x$-axis, which is to be denoted by $N^{**}$ in accordance with the notation we have used throughout, represents the image of the point $N^0$, and that the ratio of the linear dimensions of the parts of a composite image to those of the object is constant, namely $g$ to $1$ or $1$ to $l$.

\Article{12}
The case considered in the preceding article occurs in a telescope whose lenses are adjusted for a farsighted eye and for a clear view of objects at an infinite distance. The formulas above show that the direction of the outgoing ray depends only on the direction of the incident ray. Thus incident rays parallel to one another correspond to outgoing rays that are also parallel, and the tangent of the angle the incident rays make with the axis is to the tangent of the angle the outgoing rays make with it as $1$ is to $l$. The number $l=\frac1g$ is therefore what is called the magnification of the telescope. Its positive or negative sign indicates whether the objects appear upright or inverted. If the incident and outgoing rays are made to exchange roles by turning the eyepiece side toward the objects. the objects appear reduced in the same ratio, and this is the basis of the method for determining the magnification of a tele\SourceLeaf{267}scope, as convenient as it is precise, which I reported in 1823 in Volume 2 of the \textit{Astronomischen Nachrichten}.

Another method of determining the magnification compares an object with its image by the ratio of their linear dimensions. \textsc{Ramsden}'s dynameter is simply a device for measuring the diameter of the image, formed at $N^{**}$, of the circular boundary of the objective. One must of course first make sure that this image actually appears and is not obscured by an internal stop. The image must also be real, which requires $ghn^*$ to be negative. In a \textsc{Galilei}an telescope, where this image is only virtual, an accurate result could be obtained only with a micrometric microscope. That microscope would also deserve preference in \textit{all} cases. where greater precision is desired. The preceding article also shows that a suitable object at some distance from the objective can be used just as well, provided only that the distance does not become so great that the image ceases to be real or within reach of the microscope. Finally, it may be noted that $N^{**}$ is the point known in the theory of telescopes as the \textit{position of the eye}.

\Article{13}
To apply the general rules of Article 2 to a simple glass lens, we denote the ratio of refraction in passing from air into glass by $n:1$, the radii of the first and second surfaces by $(n-1)f$ and $(n-1)f'$, and the thickness of the lens by $ne$. In place of the notation there, we therefore have
\[
\begin{array}{ccc}
n^0&\mathit{here}&1\\
n'&.\,.\,.\,.&n\\
n''&.\,.\,.\,.&1\\
t'&.\,.\,.\,.&e\\
u^0&.\,.\,.\,.&-\dfrac1f\\\noalign{\vskip 4pt}
u'&.\,.\,.\,.&-\dfrac1{f'}
\end{array}
\]
and consequently
\SourceLeaf{268}
\[
\begin{aligned}
g&=1+u^0t'=\frac{f-e}{f}\\
h&=t'=e\\
k&=u^0+u'+t'u^0u'=-\frac{f+f'-e}{ff'}\\
l&=1+u't'=\frac{f'-e}{f'}
\end{aligned}
\]
Thus for the focal length $\varphi$ we have, by Article 9,
\[
\varphi=\frac{ff'}{f+f'-e}
\]
for the two principal points, denoted here by $E,E'$, by Article 6,
\[
\begin{aligned}
E&=N^0+\frac{ef}{f+f'-e}=N^0+\frac{e\varphi}{f'}\\
E'&=N'-\frac{ef'}{f+f'-e}=N'-\frac{e\varphi}{f}
\end{aligned}
\]
and for the two focal points $F,F'$
\[
\begin{aligned}
F&=E-\varphi=N^0-\frac{f(f'-e)}{f+f'-e}\\
F'&=E'+\varphi=N'+\frac{f'(f-e)}{f+f'-e}
\end{aligned}
\]
For the intersection with the axis of the straight line traced by a principal ray inside the lens, extended forward or backward if necessary, one readily finds
\[
x=N^0+\frac{nef}{f+f'}=N'-\frac{nef'}{f+f'}
\]
Some writers call this point, which is therefore independent of the inclination of the principal ray, the optical center of the lens. It scarcely seems to deserve this distinction, since it has no other noteworthy properties. The name seems to have led here and there to the mistaken belief that the simple relations between image and object that hold for an infinitely thin lens can be transferred to a lens of finite thickness simply by using that center as the point of reference. Yet, as shown above, this transfer is valid only when the object is referred to the first principal point and the image to the second. In a system of several lenses, even in an achromatic doublet objective, there can in any case be no question of a center in that sense. If the name is to be retained, I would think it more appropriate to apply it to the point midway between the two
\SourceLeaf{269}principal points (and hence also midway between the two focal points), which coincides with the former point only if the lens is symmetrical. This point has the practical advantage that it can be located easily and precisely by reversing the lens. It is evidently the point that must return to its former position when the lens is reversed if the location of the image of a fixed object is to stay the same.

It may also be noted here that the distance between the two principal points is
\[
E'-E=ne-\frac{e(f+f')}{f+f'-e}=(n-1)e-\frac{ee}{f+f'-e}
\]
and so, insofar as $e$ is usually very small compared with $f+f'-e$, it differs scarcely at all from $(n-1)e$, or from the thickness of the lens multiplied by $\frac{n-1}{n}$.

\Article{14}
For a system of lenses on a common axis, more convenient formulas can replace the general formulas of Article 2 that determine the path of the outgoing ray from that of the incident light ray. To do so, one replaces the radii of the individual refracting surfaces and the distances between them with the focal lengths of the individual lenses and the distances from their second principal points to the first principal points of the following lenses. The new formulas are very similar to those of Article 2, but contain only half as many elements. Since they are very easy to derive from what precedes, it will be enough to give them here in a form ready for use.

We denote the focal lengths of the successive individual lenses by $\varphi^0,\varphi',\varphi''$, and so on. For their principal points we depart here from the notation used so far. We denote the first principal points by $E^0,E',E''$, and so on, and the second by $I^0,I',I''$, and so on. For brevity we write
\[
\begin{aligned}
-\frac1{\varphi^0}&=u^0,&-\frac1{\varphi'}&=u',&-\frac1{\varphi''}&=u''\quad\text{etc.}\\
E'-I^0&=t',&E''-I'&=t'',&E'''-I''&=t'''\quad\text{etc.}
\end{aligned}
\]
The last terms in these series may be marked as such by an asterisk.

\SourceLeaf{270}If one now writes the equations for the incident ray in the form
\[
\begin{aligned}
y&=\beta^0(x-E^0)+b^0\\
z&=\gamma^0(x-E^0)+c^0
\end{aligned}
\]
and those for the outgoing ray in the following form
\[
\begin{aligned}
y&=\beta^*(x-I^*)+b^*\\
z&=\gamma^*(x-I^*)+c^*
\end{aligned}
\]
then, if the four quantities $g,h,k,l$ are determined by formulas identical to those designated (5) in Article 2, one obtains
\[
\begin{aligned}
b^*&=gb^0+h\beta^0,&c^*&=gc^0+h\gamma^0\\
\beta^*&=kb^0+l\beta^0,&\gamma^*&=kc^0+l\gamma^0
\end{aligned}
\]
For the two principal points of the lens system considered as a whole, the coordinates are
\[
\begin{aligned}
\text{for the first}\quad x&=E^0-\frac{1-l}{k}\\
\text{for the second}\quad x&=I^*+\frac{1-g}{k}
\end{aligned}
\]
The coordinates of the two focal points of the lens system are
\[
\begin{aligned}
\text{for the first}\quad x&=E^0+\frac{l}{k}\\
\text{for the second}\quad x&=I^*-\frac{g}{k}
\end{aligned}
\]
\[
\text{the focal length itself is}\quad=-\frac1k
\]

The formulas for the case where the system consists of only two lenses deserve to be written out separately. They are
\[
\begin{aligned}
g&=\frac{\varphi^0-t'}{\varphi^0}\\
h&=t'\\
k&=-\frac{\varphi^0+\varphi'-t'}{\varphi^0\varphi'}\\
l&=\frac{\varphi'-t'}{\varphi'}
\end{aligned}
\]
The values of $x$ for the two principal points are
\SourceLeaf{271}
\[
E^0+\frac{t'\varphi^0}{\varphi^0+\varphi'-t'}\quad\text{and}\quad I'-\frac{t'\varphi'}{\varphi^0+\varphi'-t'}
\]
and the focal length
\[
=\frac{\varphi^0\varphi'}{\varphi^0+\varphi'-t'}
\]
One sees that these formulas are entirely analogous to those given in Article 13 for determining the principal points and focal length of a simple lens, with the quantities $f^0,f',e$ used there replaced here by $\varphi^0,\varphi',t'$.

The distance between the two principal points in the case of two lenses is
\[
\begin{aligned}
&=I'-E^0-\frac{t'(\varphi^0+\varphi')}{\varphi^0+\varphi'-t'}\\
&=I^0-E^0+I'-E'-\frac{t't'}{\varphi^0+\varphi'-t'}
\end{aligned}
\]
If $t'$ is very small, as is always the case with achromatic doublets of the usual construction, the last term becomes insignificant. The distance between the two principal points for such a doublet considered as a whole is therefore very nearly equal to the sum of the two values that this distance has in the individual lenses.

It is also self-evident that all the formulas given in this article can be applied without any change when partial systems of lenses are to be combined into one whole system in place of simple lenses.

\Article{15}
As we have shown, the optical phenomena produced by a single lens or by a system of several lenses on a common axis depend on three elements. These are determined by the ratio of refraction (or by the ratios of refraction if they differ among the lenses) and by the positions and radii of the refracting surfaces. These quantities are usually not known directly, however, so something remains to be said about the converse method of deriving the three elements from observed phenomena. We denote the various points on the axis involved here as follows:
\SourceLeaf{272}$\xi$ is an object and $\xi'$ its image. $F$ and $F'$ are the first and second focal points, $E$ and $E'$ the first and second principal points, and finally $D$ is a point fixed relative to the lens (or lens system). As always, the same letters denote the coordinates of these points in each experiment. We further set the focal length $=f$ and write the distances of the point $D$ from the focal points as $D-F=p$ and $F'-D=q$. The three quantities $f,p,q$ can be regarded as the elements of the lens. Three experiments will therefore always be needed to determine them. The distances of the object and its image from the point $D$ must be measured for three different positions of the object and image relative to the lens. We shall first solve this problem in full generality.

Let the values of $D-\xi$ and $\xi'-D$ be $a,b$ in one experiment, $a',b'$ in a second and $a'',b''$ in a third. The general equation
\[
(F-\xi)(\xi'-F')=ff
\]
therefore gives us
\[
\begin{aligned}
(a-p)(b-q)&=ff\\
(a'-p)(b'-q)&=ff\\
(a''-p)(b''-q)&=ff
\end{aligned}
\]
from which elimination readily yields
\[
\begin{aligned}
p&=a-\frac{(a'-a)(a''-a)(b'-b'')}{R}\\
q&=b-\frac{(b-b')(b-b'')(a''-a')}{R}\\
ff&=\frac{(a'-a)(a''-a)(a''-a')(b-b')(b-b'')(b'-b'')}{RR}
\end{aligned}
\]
where, for brevity, we write
\[
(a''-a)(b-b')-(a'-a)(b-b'')=R
\]
The quantity $R$ can also be put in the following form
\[
\begin{aligned}
R&=(a''-a')(b-b')-(a'-a)(b'-b'')\\
&=(a''-a')(b-b'')-(a''-a)(b'-b'')
\end{aligned}
\]
and $p$ and $q$ in the following forms
\SourceLeaf{273}
\[
\begin{aligned}
p&=a'-\frac{(a'-a)(a''-a')(b-b'')}{R}\\
&=a''-\frac{(a''-a)(a''-a')(b-b')}{R}\\
q&=b'-\frac{(b-b')(b'-b'')(a''-a)}{R}\\
&=b''-\frac{(b-b'')(b'-b'')(a'-a)}{R}
\end{aligned}
\]

\Article{16}
A few remarks must be added to the general solution given in the preceding article.

I. The object is assumed to lie on the same side of the lens in all three experiments. If it is convenient to reverse the lens in one of the experiments, that experiment can be regarded as though the image were the object and the object the image. This reduces the case to the preceding one.

II. Taken by itself, the formula for $ff$ leaves open whether $f$ is to be taken as positive or negative. This is already decided by whether the image is upright or inverted, since $\xi'-F'$ and $f$ must have opposite signs in the first case and the same sign in the second. It should also be noted that, for all the generality of the analytical solution, its practical application is limited to real images (and thus, for individual lenses, to positive focal lengths), unless special means are brought in to determine the position of virtual images.

III. Since the experiments can be performed only to a limited degree of precision, the reliability of the results depends very much on which combinations are chosen. As a general rule, three experiments, two of them carried out under conditions that differ little from each other, cannot determine all three elements precisely.

\Article{17}
In a simple lens and likewise in one made up of two or more lenses lying very close together (as in achromatic objectives of the usual construction), the two principal \SourceLeaf{274}points lie a short distance apart. If this distance $E'-E=\lambda$ could be treated as a known quantity, \textit{two} experiments would suffice, since the equation
\[
p+q=2f+\lambda
\]
takes the place of the third experiment. Combining it with the other two equations
\[
\begin{aligned}
(a-p)(b-q)&=ff\\
(a'-p)(b'-q)&=ff
\end{aligned}
\]
and eliminating $p$ and $q$ gives the equation for determining $f$
\[
\begin{aligned}
\frac{(a'+b'-a-b)^2}{(a'-a)(b-b')}.ff+2(a+b+a'+b'-2\lambda)f\\
{}-(a+b'-\lambda)(a'+b-\lambda)=0
\end{aligned}
\]
This quadratic equation becomes linear when $a'+b'-a-b=0$, that is, when the two experiments are arranged so that the distance of the image from the object stays the same in both while the lens occupies two different positions. Let this distance be $=c$, so that $a=c-b$ and $a'=c-b'$. This gives
\[
4(c-\lambda)f=(c-\lambda+b'-b)(c-\lambda-b'+b)
\]
or
\[
f=\tfrac14(c-\lambda)-\frac{(b'-b)^2}{4(c-\lambda)}
\]
For each prescribed value of $c$, $F-\xi$ must satisfy the following equation.
\[
F-\xi+\frac{ff}{F-\xi}=F-\xi+\xi'-F'=c-2f-\lambda
\]
Its two roots
\[
\begin{aligned}
F-\xi&=\tfrac12(c-2f-\lambda)+\tfrac12\BareRoot{(c-4f-\lambda)(c-\lambda)}\\
F-\xi&=\tfrac12(c-2f-\lambda)-\tfrac12\BareRoot{(c-4f-\lambda)(c-\lambda)}
\end{aligned}
\]
are real and unequal when $c$ is greater than $4f+\lambda$, so that for a fixed object $\xi$ there are then always two different positions of the lens for which the image coincides with the point $\xi+c$. The product of these two values of $F-\xi$, that is, $(a-p)(a'-p)$, is $=ff$, from which it is also clear
\SourceLeaf{275}that $a'-p=b-q$ and $b'-q=a-p$, and consequently
\[
\begin{aligned}
p&=\tfrac12(2f+c+\lambda-b-b')\\
q&=\tfrac12(2f-c+\lambda+b+b')\\
E&=D+\tfrac12(b+b'-c)-\tfrac12\lambda\\
E'&=D+\tfrac12(b+b'-c)+\tfrac12\lambda
\end{aligned}
\]

\Article{18}
At the position of the lens where $F-\xi=f$, $\xi'-\xi=4f+\lambda$ holds, or the image is at its minimum distance from the object. It moves away from the object as soon as the lens is shifted from that position to either side, though evidently very slowly at first. It follows that if the value chosen for $c$ exceeds $4f+\lambda$ only slightly, the experiments for determining the two required positions of the lens, or the values of $b$ and $b'$, allow only comparatively low precision. This uncertainty carries over in full to the determination of $E$ and $E'$. Under such circumstances the method is therefore scarcely usable for \textit{this} purpose. The situation is different if the aim is only to determine the focal length. That circumstance causes no loss of precision, since the expression for $f$ contains only the square of $b'-b$. The method is also easier to carry out in this case, since in addition to the distance $c$ only the amount by which the lens is shifted, $b'-b$, need be measured. The absolute values of $b$ and $b'$ are therefore unnecessary.

\Article{19}
If one neglects $\lambda$ entirely and thus sets
\[
f=\tfrac14c-\frac{(b'-b)^2}{4c}
\]
the method described agrees with the one that \textsc{Bessel} proposed in Volume 17 of the Astronomischen Nachrichten and applied to determining the focal length of the objective of the Königsberg heliometer. The exact formula shows that when $\lambda$ is neglected, the focal length obtained is too large by
\SourceLeaf{276}
\[
\tfrac14\lambda+\frac{\lambda(b'-b)^2}{4c(c-\lambda)}
\]
where the second term may be regarded as imperceptible under the circumstances mentioned. To obtain a result commensurate with the precision that the method itself allows, it is therefore essential to take $\lambda$ into account. Obtaining an accurate knowledge of this quantity presents some difficulty, however. For a simple lens it will suffice to calculate the approximate value for $\lambda$ given above in Article 13 from the measured thickness and the roughly known ratio of refraction. For an achromatic doublet one might also use the approximate value given above in Article 14, provided one can determine the thickness of each component accurately. To get at least a rough idea of the effect that neglecting $\lambda$ can have, let us consider, by way of example, an objective whose crown glass lens is 7 lines thick and whose flint glass lens is 3 lines thick. Let us take the ratios of refraction to be 1,528 for the former and 1,618 for the latter. This makes the separation of the two principal points approximately
\[
\begin{array}{lr}
\text{for the crown glass lens}&2{,}42\\
\text{for the flint glass lens}&1{,}15\\
\text{and for the combined lens}&3{,}57\ \text{lines}
\end{array}
\]
so that the focal length is too large by 0,89 lines. For an objective with a focal length of 8 feet and the dimensions assumed here, the error would therefore amount to about $\frac1{1300}$ of the whole.

\Article{20}
If the method of determining $\lambda$ given in the preceding article cannot be used, or if one does not wish to settle for it, the following approach seems the most suitable way to find it by direct experiments.

Determine the position of the image of a very distant object (as nearly on the axis of the lens as possible) relative to the point $D$. Insofar as the distance of the object can be regarded as infinitely large, the image falls
\SourceLeaf{277}at $F'$, and the measured distance $F'-D$ gives $q$ directly. Repeat the experiment with the lens reversed. The image will then fall at $F$, and its distance from $D$ will give the value of $p$. For the third experiment, bring the object (on the side of $F$) as close to the lens as possible. Determine the distance of the image from this object $=\xi'-\xi$ and at the same time the distance $D-\xi=a''$, and set $\xi'-D=\xi'-\xi-a''=b''$. This gives
\[
(p-a'')(q-b'')=ff
\]
or
\[
\lambda=p+q-2\BareRoot{(p-a'')(q-b'')}
\]
If the measurements in all three experiments could be made with the greatest precision, they alone suffice to determine all three elements $p,q,f$ accurately enough, and no further experiments are needed. If one wishes to obtain $f$ with still greater precision, these experiments should be regarded only as preparation for the procedure of Article 18. That preparation provides the value of $\lambda$. To see more clearly what chiefly governs the precision with which $\lambda$ is determined in this way, we put the above formula for $\lambda$ in the following form
\[
\lambda=a''+b''+\frac{(p-a''-\BareRoot{(p-a'')(q-b'')})^2}{p-a''}
\]
and consider that $p-a''-\BareRoot{(p-a'')(q-b'')}$ represents the distance of the object in the third experiment from the first principal point, while $p-a''$ represents the distance of that object from the first focal point. It is clear from this that under the conditions in question, the last term of the formula for $\lambda$ becomes very small and that small inaccuracies in the values of $p,q,a'',b''$ have little effect on its calculated value. The precision of the determination of $\lambda$ therefore depends chiefly on the precision with which $\xi'-\xi=a''+b''$ is measured.

\Article{21}
A few remarks on the procedure described in the preceding article deserve a place here.

I. The means ordinarily available are insufficient for carrying out the third experiment, where the image is only virtual. The following method, however, com\SourceLeaf{278}bines convenience and precision. Draw a circle on a flat surface, as large as or slightly larger than the projecting rim of the lens mount, and mark its center with two fine intersecting lines. Place the lens and its mount on the surface so that the mount is concentric with the circle. Then set a compound microscope fitted with a crosshair vertically above it on a fixed stand and slide the microscope in its sleeve until the image of the cross mark coincides exactly with the crosshair. Finally remove the lens and bring the microscope closer to the plane by sliding it in its sleeve until the image of the cross mark again coincides with the microscope's crosshair. The amount of this last displacement is easy to measure precisely by any suitable means. It is the distance of the object (the cross mark) from its image produced by the glass line $=\xi'-\xi$. One may take the point on the axis of the lens that lies in the plane touching the projecting rim of the mount as the fixed point $D$ itself. In that case $a''=0$ and $b''=\xi'-\xi$. If a different point $D$ was chosen, one can compare the two points by readily available means to find $a''$.

II. If the object used for the first and second experiments is far away, yet its image cannot be regarded as coinciding exactly with the focal point, a correction is needed. It is obtained by dividing the square of the focal length by the distance of the object. Subtract this correction from the distances of the image from the point $D$ to obtain the quantities $q$ and $p$ exactly. Clearly, only a rough approximation to the focal length and the distance is needed for this, insofar as the distance is very large. The correction can, however, be determined just as easily by direct, rigorously exact formulas. For the first experiment let $a$ be the value of $D-\xi$ and $b$ the value of $\xi'-D$. For the second experiment, in which the lens is used in the reverse orientation, denote the distance $D-\xi$ by $b'$ and $\xi'-D$ by $a'$. This choice, which amounts to the same thing as the one given in Article 16, I, gives us the advantage that the equations for the three experiments
\SourceLeaf{279}
\[
\begin{aligned}
(a-p)(b-q)&=ff\\
(b'-q)(a'-p)&=ff\\
(a''-p)(b''-q)&=ff
\end{aligned}
\]
are identical to those from which we started in Article 15. The formulas derived from them by elimination therefore still hold. If the second experiment was carried out in such a way that the position of the image in space is the same as in the first (which is easily done, although at the assumed large distance of the object the result is unaffected to any perceptible degree even if one is less exacting about this), then $a+b=b'+a'$. We denote this quantity by $c$, which further simplifies the formulas of Article 15. Namely, the second formula for $p$ and the first for $q$ give
\[
\begin{aligned}
p&=a'-\frac{(a'-a'')(b-b'')}{c-a''-b''}\\
q&=b-\frac{(a'-a'')(b-b'')}{c-a''-b''}
\end{aligned}
\]

III. Even if one wishes to reserve the most precise determination of the focal length for the method of Article 17 instead of using the procedure of Article 20 to determine all the elements, the latter procedure is still the most suitable for establishing the positions of the two principal points. Namely,
\[
\begin{aligned}
E&=D+\tfrac12(q-p)-\tfrac12\lambda=D+\tfrac12(b-a')-\tfrac12\lambda\\
E'&=D+\tfrac12(q-p)+\tfrac12\lambda=D+\tfrac12(b-a')+\tfrac12\lambda
\end{aligned}
\]

\Article{22}
For a simple lens, and more generally for a system of lenses, a definite value can be assigned to the focal length and definite positions to the principal and focal points only when the light rays have a definite refrangibility. For rays of a different refrangibility those points take other positions and the focal length takes another value. The nonhomogeneous light from objects therefore undergoes chromatic dispersion as it passes through glass. This dispersion can be eliminated by combining two or
\SourceLeaf{280}more lenses made of different kinds of glass. For an achromatic objective to be perfect, however, parallel rays must converge at one point regardless of their color. This applies both to rays parallel to the axis and to those incident at an angle to it. In other words, the differently colored images of an extended object regarded as infinitely distant must lie in one plane and must also be the same size. The first condition rests on the second focal point being the same for differently colored rays, and the second on equality of focal length. Since the focal length is the distance of the second focal point from the second principal point, one can also express both conditions by saying that both points must be the same for red and violet rays alike. If only the first condition is satisfied, rays inclined to the axis do not give a pure image. A \textit{very} slight difference in focal length for differently colored rays may always be regarded as quite harmless, however.

In the theory of achromatic objectives, it is customary to take only the first condition into account. Yet in the usual construction of these objectives, where the two lenses are either in contact or separated by an extremely small distance, the different refrangibilities of the light rays have so little effect on the positions of the principal points that the second condition is satisfied of itself, either exactly or so nearly that no perceptible imperfection can result. If one considers it worth the trouble, one can also calculate the thicknesses of the lenses so that the second principal point is exactly the same for rays of different refrangibility.

The situation is different if the convex crown glass lens is separated from the concave flint glass lens by a considerable distance. It is easy to show that for choices of this distance and of the focal lengths of the individual lenses that keep the second focal point of the lens system the same for differently colored rays, the focal length of the system necessarily becomes greater for violet rays than for red ones. The difference is of the same order as that found, in the opposite direction, in simple lenses. The same still holds if, as in so-called dialytic telescopes, the second lens is replaced by
\SourceLeaf{281}a combination of a flint glass lens and a crown glass lens, either in contact or very slightly separated. It always remains impossible in this way to produce an image of an extended object entirely free of color errors. If the violet image is to be the same distance from the lens system as the red one, it necessarily becomes larger than the red one.

One must by no means conclude from this that \textit{telescopes} of the latter design must remain more imperfect in respect of achromatism than telescopes with achromatic objectives of the usual construction that produce an image entirely free of color errors. On the contrary, one can claim the very reverse. With a carefully calculated arrangement of the eyepieces, telescopes of the latter design can present the eye with an image with less color error.

Indeed, to the \textit{eye}, even a perfectly pure image produced by the objective, whether real or virtual, cannot \textit{appear} perfectly pure because of the dispersion produced by the eyepiece lenses. A special arrangement of the eyepieces prevents the so-called colored fringe, but cannot eliminate the longitudinal deviation, which is further increased by the fact that the human eye itself is not achromatic. This arrangement makes the final images, red and violet, appear the same size. It does not make them appear at the same distance or sharply \textit{at the same time}.

The unequal sizes of the first images, red and violet, are unavoidable with dialytic objectives. They can, however, be compensated quite well by a suitable design of the eyepieces, so that the visible colored fringe is eliminated just as well as in telescopes of the usual design. The second imperfection just mentioned persists here too as long as the first red and violet images lie at the same distance from the objective.

It is therefore clear that to produce an image entirely free of color errors \textit{in the eye}, the first image \textit{must} have a certain longitudinal deviation that depends on the properties of the eyepieces and on the lack of achromatism of the human eye. In theory, an objective of the usual design can indeed also be calculated to produce a prescribed longitudinal deviation. Yet even apart from the difficulty of achieving in manufacture all the
\SourceLeaf{282}precision required to realize such very small differences, this longitudinal deviation would still suit only one particular eyepiece. In the dialytic design, however, the doublet that forms the second part of the objective can be moved relative to the first. This makes it possible to obtain the longitudinal deviation required for each eyepiece, while the eyepiece can be designed to eliminate the colored fringe as required. For the rest, I must confine myself here to this brief indication and leave a fuller treatment of this interesting subject for another occasion.

